\documentclass[12pt]{article}

\usepackage[portuguese]{babel}
\usepackage{float}
\usepackage{array}
\usepackage{listings}
\usepackage{indentfirst}
\usepackage{amsfonts, amssymb, amsmath}
\usepackage{graphicx}
\usepackage{enumitem}
\usepackage[colorlinks=true, allcolors=blue]{hyperref}
\usepackage[letterpaper,
            top=1.5cm,
            bottom=1.5cm,
            left=2.75cm,
            right=2.75cm,
            marginparwidth=1.75cm]{geometry}

\title{Arquitetura e Organização de Computadores II}
\author{Prof. Anderson}
\date{\today}

\begin{document}

\maketitle

\begin{itemize}
    \item Arquiteturas paralelas
    \item Topologia
    \item Coerência da Cache
    \item Multinucleo
    \item Paralelismo -- Dados
    \item GPU
    \item Cluster
\end{itemize}

Quando um processador é escalar, isso indica que ele faz uma operação por 
unidade de tempo

Em Arq I vemos processadores escalares, em Arq II veremos processadores 
superescalares

\section{Paralelismo e Desempenho}


\subsection{Desempenho}
\begin{itemize}
    \item Speedup -- \( S = \frac{T_s}{T_p}\)
    \begin{itemize}
        \item $T_s$ -- Tempo sequencial
        \item $T_p$ -- Tempo paralelo
    \end{itemize}
\end{itemize}

Trabalho 1 ou 11

Polybench
cBench
NAS

rodinia

\end{document}